% Part:sets-functions-relations
% Chapter: size-of-sets
% Section: non-enumerability

\documentclass[../../../include/open-logic-section]{subfiles}

\begin{document}

\olfileid{sfr}{siz}{nen}

\olsection{\printtoken{S}{nonenumerable} সেট}

\begin{editorial}
  এই অংশে \olref[enm]{sec}-এর সংজ্ঞা ব্যবহার করে $\Bin^\omega$ ও $\Pow{\PosInt}$-এর অতালিকায়নযোগ্যতা প্রমাণ করা হয়েছে। এটি \olref[enm-alt]{sec}-এর সংস্করণের তুলনায় কিছুটা বেশি প্রাথমিক ও বিশদভাবে লেখা।
\end{editorial}

ধনাত্মক পূর্ণসংখ্যার সেট $\PosInt$-এর মতো কিছু সেট অসীম। এ পর্যন্ত আমরা অসীম সেটের যেসব উদাহরণ দেখেছি, সেগুলি সবই !!{enumerable}। কিন্তু এমন অসীম সেটও আছে, যাদের এই ধর্ম নেই। এই ধরনের সেটকে \emph{!!{nonenumerable}} বলা হয়।

প্রথমেই !!{nonenumerable} সেটের অস্তিত্ব হয়তো বিস্ময়কর মনে হবে। যেকোনো !!{enumerable} সেট~$A$-এর জন্য !!a{surjective} অপেক্ষক $f \colon \PosInt \to A$ আছে। কোনো সেট !!{nonenumerable} হলে এমন অপেক্ষক নেই। অর্থাৎ~$\PosInt$-এর অসীমসংখ্যক !!{element}s-কে~$A$-তে পাঠানো কোনো অপেক্ষকই~$A$-এর সবটুকু নিঃশেষ করতে পারে না। তাই~$A$-তে অসীমসংখ্যক ধনাত্মক পূর্ণসংখ্যার চেয়েও ``বেশি'' !!{element}s আছে।

একটি সেট !!{nonenumerable}, তা কীভাবে প্রমাণ করা যায়? এমন কোনো সমাপতিত অপেক্ষকের অস্তিত্ব থাকতে পারে না, তা দেখাতে হবে। সমতুল্যভাবে দেখাতে হবে যে~$A$-এর উপাদানগুলিকে একদিকে অসীম তালিকায় তালিকায়ন করা যায় না। এর সবচেয়ে ভালো উপায় হলো দেখানো যে~$A$-এর !!{element}s-এর প্রত্যেক তালিকাতেই অন্তত একটি উপাদান বাদ পড়বে; অথবা কোনো অপেক্ষক $f\colon \PosInt \to A$ !!{surjective} হতে পারে না। কান্টরের \emph{কর্ণ পদ্ধতি} ব্যবহার করে এটি করা যায়।~$A$-এর !!{element}s-এর কোনো তালিকা, ধরা যাক $x_1$, $x_2$, \dots, দেওয়া থাকলে আমরা~$A$-এর আরেকটি উপাদান গঠন করি, যা গঠনপ্রণালির কারণেই ওই তালিকায় থাকতে পারে না।

আমাদের প্রথম উদাহরণ হলো $0$ ও $1$-এর সকল অসীম, ফাঁকবিহীন অনুক্রমের সেট~$\Bin^\omega$।

\begin{thm}
\ollabel{thm:nonenum-bin-omega}
$\Bin^\omega$~!!{nonenumerable}।
\end{thm}

\begin{proof}
স্ববিরোধে পৌঁছানোর উদ্দেশ্যে ধরি, $\Bin^\omega$ !!{enumerable}, অর্থাৎ ধরি $\Bin^\omega$-এর সকল !!{element}s-এর একটি তালিকা $s_{1}$, $s_{2}$, $s_{3}$, $s_{4}$, \dots{} আছে। এই $s_i$-গুলির প্রত্যেকটিই $0$ ও~$1$-এর একটি অসীম অনুক্রম। এই তালিকার $i$-তম অনুক্রমের $j$-তম উপাদানকে $s_i(j)$ বলি। তাহলে $i$-তম অনুক্রম~$s_i$ হলো
\[
s_i(1), s_i(2), s_i(3), \dots
\]

এই তালিকা এবং এর প্রত্যেক অনুক্রম $s_i$-এর উপাদানগুলিকে একটি সারণিতে সাজাতে পারি:
\[
\begin{array}{c|c|c|c|c|c}
& 1 & 2 & 3 & 4 & \dots \\\hline
1 & \mathbf{s_{1}(1)} & s_{1}(2) & s_{1}(3) & s_1(4) & \dots \\\hline
2 & s_{2}(1)& \mathbf{s_{2}(2)} & s_2(3) & s_2(4) & \dots \\\hline
3 & s_{3}(1)& s_{3}(2) & \mathbf{s_3(3)} & s_3(4) & \dots \\\hline
4 & s_{4}(1)& s_{4}(2) & s_4(3) & \mathbf{s_4(4)} & \dots \\\hline
\vdots & \vdots & \vdots & \vdots & \vdots & \mathbf{\ddots}
\end{array}
\]
বাঁ দিকের চিহ্নগুলি তালিকার অনুক্রম $s_1$, $s_2$, \dots-এর সংখ্যা দেয়; উপর দিকের সংখ্যাগুলি প্রতিটি অনুক্রমের !!{element}s চিহ্নিত করে। যেমন, $s_{1}(1)$ হলো অনুক্রম $s_{1}$-এর প্রথম !!{element}, সেটি $0$ বা~$1$ যে সংখ্যাই হোক তার একটি নাম; অন্যগুলির ক্ষেত্রেও একই কথা।

এখন $0$ ও $1$-এর এমন একটি অসীম অনুক্রম $\overline{s}$ গঠন করব, যা এই তালিকায় কোনোভাবেই থাকতে পারে না। $\overline{s}$-এর সংজ্ঞা তালিকা $s_1$, $s_2$, \dots-এর উপর নির্ভর করবে। $0$ ও $1$-এর অসীম অনুক্রমগুলির যেকোনো অসীম তালিকা এমন একটি অসীম অনুক্রম~$\overline{s}$ নির্ধারণ করে, যা নিশ্চিতভাবেই তালিকাটিতে থাকবে না।

$\overline{s}$ সংজ্ঞায়িত করতে তার সব !!{element}s নির্দিষ্ট করি, অর্থাৎ সকল $n \in \PosInt$-এর জন্য $\overline{s}(n)$ নির্দিষ্ট করি। উপরের সারণির কর্ণ বরাবর নিচে পড়ে (এই কারণেই নাম ``কর্ণ পদ্ধতি'') প্রত্যেক $1$-কে $0$ এবং প্রত্যেক $0$-কে~$1$ করে এটি করি। আরও বিমূর্তভাবে, কর্ণের $n$-তম !!{element} $s_n(n)$ যথাক্রমে $1$ বা $0$ হলে $\overline{s}(n)$-কে $0$ বা $1$ সংজ্ঞায়িত করি।
\[
\overline{s}(n) =
\begin{cases}
1 & \text{যদি $s_{n}(n) = 0$ হয়}\\
0 & \text{যদি $s_{n}(n) = 1$ হয়}.
\end{cases}
\]
ক্ষেত্রভিত্তিক সংজ্ঞার চেয়ে সূত্র পছন্দ হলে $\overline{s}(n) = 1 - s_n(n)$ দিয়েও সংজ্ঞায়িত করতে পারো।

স্পষ্টতই $\overline{s}$ হলো $0$ ও $1$-এর একটি অসীম অনুক্রম, কারণ আমাদের সারণির কর্ণে থাকা $0$ ও $1$-এর অনুক্রমটির প্রত্যেক বিট উল্টে এটি পাওয়া যায়। তাই $\overline{s}$ হলো~$\Bin^\omega$-এর !!a{element}। কিন্তু এটি তালিকা $s_1$, $s_2$, \dots{}-এ থাকতে পারে না। কেন?

এটি তালিকার প্রথম অনুক্রম $s_1$ হতে পারে না, কারণ প্রথম !!{element}-এ এটি $s_1$-এর থেকে ভিন্ন। $s_1(1)$ যা-ই হোক, আমরা $\overline{s}(1)$-কে তার বিপরীত মান হিসেবে সংজ্ঞায়িত করেছি। এটি তালিকার দ্বিতীয় অনুক্রমও হতে পারে না, কারণ দ্বিতীয় উপাদানে $\overline{s}$ ও $s_2$ ভিন্ন: $s_2(2)$ যদি $0$ হয়, তবে $\overline{s}(2)$ হলো $1$, এবং উল্টোটিও সত্য। এইভাবে চলতে থাকে।

আরও নির্দিষ্টভাবে: $\overline{s}$ তালিকায় থাকলে এমন কোনো $k$ থাকত, যাতে $\overline{s} = s_{k}$। দুটি অনুক্রম অভিন্ন হয় যদি এবং কেবল যদি তারা প্রত্যেক স্থানে একমত হয়, অর্থাৎ যেকোনো~$n$-এর জন্য $\overline{s}(n) = s_{k}(n)$। তাই বিশেষ ক্ষেত্রে $n = k$ নিয়ে $\overline{s}(k) = s_{k}(k)$ হতে হতো। $s_k(k)$ হয় $0$, নয়তো~$1$। এটি $0$ হলে $\overline{s}(k)$ অবশ্যই~$1$—কারণ আমরা এভাবেই $\overline{s}$ সংজ্ঞায়িত করেছি। আবার $s_k(k) = 1$ হলে, আবারও $\overline{s}$-এর সংজ্ঞার কারণে $\overline{s}(k) = 0$। উভয় ক্ষেত্রেই $\overline{s}(k) \neq s_{k}(k)$।

শুরুতে আমরা ধরে নিয়েছিলাম যে $\Bin^\omega$-এর !!{element}s-এর একটি তালিকা $s_1$, $s_2$, \dots{} আছে। এই তালিকা থেকে আমরা এমন একটি অনুক্রম~$\overline{s}$ গঠন করেছি, যা তালিকায় থাকতে পারে না বলে প্রমাণ করেছি। কিন্তু সকল $s_i$ যদি $0$ ও $1$-এর অনুক্রম হয়, তবে এটি নিশ্চিতভাবেই $0$ ও $1$-এর অনুক্রম, অর্থাৎ $\overline{s} \in \Bin^\omega$। বিশেষত এর থেকে বোঝা যায় যে~$\Bin^\omega$-এর \emph{সকল} !!{element}s-এর কোনো তালিকা থাকতে পারে না; কারণ এমন যেকোনো তালিকা থেকে এমন একটি অনুক্রম~$\overline{s}$ গঠন করতে পারতাম, যা নিশ্চিতভাবে তালিকায় নেই। তাই~$\Bin^\omega$-এর সকল অনুক্রমের তালিকা আছে—এই অনুমান স্ববিরোধ ঘটায়।
\end{proof}

\begin{explain}
এই প্রমাণপদ্ধতিকে ``কর্ণীকরণ'' বলা হয়, কারণ এতে~$\overline{s}$ সংজ্ঞায়িত করতে সারণির কর্ণ ব্যবহার করা হয়। কর্ণীকরণে সারণি থাকতেই হবে এমন নয়: সারণি বা সত্যিকারের কর্ণ না-থাকলেও অনুরূপ ধারণা ব্যবহার করে সেট !!{enumerable} নয় দেখাতে পারি।
\end{explain}

\begin{thm}
\ollabel{thm:nonenum-pownat}
$\Pow{\PosInt}$ !!{enumerable} নয়।
\end{thm}

\begin{proof}
একইভাবে এগোব:~$\PosInt$-এর উপসেটগুলির প্রত্যেক তালিকার জন্য~$\PosInt$-এর এমন একটি উপসেট আছে, যা তালিকায় থাকতে পারে না, তা দেখাব। ধরি, নিচে~$\PosInt$-এর উপসেটগুলির একটি প্রদত্ত তালিকা আছে:
\[
Z_{1}, Z_{2}, Z_{3}, \dots
\]
এখন এমন একটি সেট $\overline{Z}$ সংজ্ঞায়িত করি, যাতে যেকোনো $n \in \PosInt$-এর জন্য $n \in \overline{Z}$ যদি এবং কেবল যদি $n \notin Z_{n}$ হয়:
\[
\overline{Z} = \Setabs{n \in \PosInt}{n \notin Z_n}
\]
$\overline{Z}$ স্পষ্টতই ধনাত্মক পূর্ণসংখ্যার একটি সেট, কারণ অনুমান অনুসারে প্রত্যেক~$Z_n$ তেমনই; তাই $\overline{Z} \in \Pow{\PosInt}$। কিন্তু $\overline{Z}$ তালিকায় থাকতে পারে না। এটি দেখাতে প্রত্যেক $k \in \PosInt$-এর জন্য $\overline{Z} \neq Z_k$ প্রতিষ্ঠা করব।

তাই যেকোনো $k \in \PosInt$ নিই। আমরা $\overline{Z}$-কে এভাবে সংজ্ঞায়িত করেছি যে, যেকোনো $n \in \PosInt$-এর জন্য $n \in \overline{Z}$ যদি এবং কেবল যদি $n \notin Z_n$। বিশেষত $n=k$ নিলে, $k \in \overline{Z}$ যদি এবং কেবল যদি $k \notin Z_k$। কিন্তু এতে বোঝা যায় $\overline{Z} \neq Z_k$, কারণ $k$ একটির !!a{element}, কিন্তু অন্যটির নয়; তাই $\overline{Z}$ ও $Z_k$-এর !!{element}s ভিন্ন। যেহেতু $k$ যেকোনো ছিল, তাই $\overline{Z}$ তালিকা $Z_1$, $Z_2$, \dots{}-এ নেই।
\end{proof}

\begin{explain}
আগের প্রমাণে কর্ণের উল্লেখ ছিল না, কিন্তু নিচের ছবির মতো ভাবলে এর মধ্যেও কর্ণ দেখতে পারো। কল্পনা করো, সেট $Z_1$, $Z_2$, \dots{} একটি সারণিতে লেখা, যেখানে প্রত্যেক !!{element}~$j \in Z_i$-কে~$j$-তম স্তম্ভে রাখা হয়েছে। ধরি, ওই তালিকার প্রথম চারটি সেট হলো $\{1,2,3,\dots\}$, $\{2, 4, 6, \dots\}$, $\{1,2,5\}$ এবং $\{3,4,5,\dots\}$। তাহলে সারণিটি এইভাবে শুরু হবে:
\[
\begin{array}{r@{}rrrrrrr}
  Z_1 = \{ & \mathbf{1}, & 2, & 3, & 4, & 5, & 6, & \dots\}\\
  Z_2 = \{ &  & \mathbf{2}, &  & 4, &  & 6, & \dots\}\\
  Z_3 = \{ & 1, & 2, &  &  & 5\phantom{,} &  & \}\\
  Z_4 = \{ &  &  & 3, & \mathbf{4}, & 5, & 6, & \dots\}\\
  \vdots & & & & & \ddots
\end{array}
\]
তখন কর্ণ বরাবর নিচে নেমে, কর্ণে থাকা সংখ্যাগুলি বাদ দিয়ে এবং $j$-তম সারি/স্তম্ভে যেখানে ফাঁক আছে সেই $j$-গুলি নিয়ে সেট $\overline{Z}$ পাওয়া যায়। উপরের ক্ষেত্রে আমরা $1$ ও $2$ বাদ দেব,~$3$ নেব,~$4$ বাদ দেব, এইভাবে চলবে।
\end{explain}

\begin{prob}
কর্ণ-যুক্তির সাহায্যে দেখাও যে $\Pow{\Nat}$ !!{nonenumerable}।
\end{prob}

\begin{prob}\label{sfr:siz:nen:prob:f-posint}
সুস্পষ্ট কর্ণ-যুক্তির সাহায্যে দেখাও যে $f \colon \PosInt \to \PosInt$ অপেক্ষকগুলির সেট !!{nonenumerable}। অর্থাৎ দেখাও, অপেক্ষকগুলির কোনো তালিকা $f_1$, $f_2$, \dots{} দেওয়া থাকলে, যেখানে প্রত্যেক $f_i\colon \PosInt \to \PosInt$, এমন কোনো $\overline{f}\colon \PosInt \to \PosInt$ আছে যা এই তালিকায় নেই।
\end{prob}

\end{document}
